\chapter*{Preface}
\label{front:preface}
\addcontentsline{toc}{chapter}{Preface}

\begin{warning}{Editorial and verification status, 4 October 2026}
\label{warn:front-verification-status}
This second draft incorporates the sequential mathematical, source,
editorial, build and rendered-page reviews of all 72 chapters and
Appendices A--N, followed by an ordered second reading for coherence
and readability. The editorial programme records the scope and evidence
of each review. Published inputs retain their stated hypotheses;
supplied manuscripts and book-level deductions are identified at their
points of use. Independent verification receipts concern their exact
dated source packets. Later editorial revisions do not acquire an earlier
receipt, and a successful compilation does not establish a theorem.
A separate foundational correction concerns the proof of homogeneous Siu
division in \cite[Theorem~6.21]{BJ22}: the changing-coefficient residual
need not decrease. \Cref{warn:ch35-siu-source-gap} records the unresolved
endpoint as $\mathrm{(SD)}_k$. The general envelope theorem and downstream
uses of that argument are conditional on this extra input; published
smooth cases with independent proofs remain available.
The current research inventory is the 42 active atoms in the twelve
problem environments of \Cref{chap:major-open-problems}. Their unresolved
hypotheses remain part of the book's mathematical boundary. In particular,
the general envelope problem over nontrivially valued fields remains open
at the scope recorded there.
\end{warning}

Non-Archimedean pluripotential theory has several origins.  One begins with
models of a projective variety over a valuation ring and regards vertical
divisors as functions on its Berkovich analytification.  A second begins with
real differential forms on tropical charts.  A third starts from filtrations
and test configurations over a trivially valued field.  A fourth starts on a
compact Kähler manifold, where geodesic rays and test curves encode objects
that behave as metrics on a non-Archimedean boundary.  Each entrance reveals
features hidden by the others.

The purpose of this book is to construct these theories carefully and to say
exactly how they are related.  The word \emph{exactly} matters.  A model psh
metric, a Chambert--Loir--Ducros psh function, a Boucksom--Favre--Jonsson psh
function, and a Boucksom--Jonsson psh function are not identical
by terminology.  They agree on important common domains because comparison
theorems make them agree.  We prove those comparison theorems wherever the
present theory allows, and we mark the hypotheses at which a comparison is
not known.

\Crefrange{part:foundations}{part:bfj} develop valued fields, models,
Berkovich spaces, tropical
differential forms, currents, and the global variational theory over a
nontrivially valued field.  \Crefrange{part:bj}{part:complex} treat the
trivially valued theory,
including the envelope theorem for unibranch projective varieties, and then
turn to complex and compact Kähler spaces.  \Cref{part:comparisons} assembles
the dictionaries among models, currents, filtrations, rays, and hybrid
degenerations.

\paragraph{How to navigate the book.}
\label{front:preface-navigation}
The introductions to the six parts give short, standard, and reference
routes, so linear reading is not required.  The global nontrivially valued
route passes from \Cref{part:foundations} to \Cref{part:bfj} and then to the
model and measure comparisons in \Cref{part:comparisons}.  The local
differential route combines the analytic-space chapters of
\Cref{part:foundations} with \Cref{part:cld}, including the
Gubler--K\"unnemann and Gubler--Rabinoff extensions.

The homogeneous and complex routes run through
\Cref{part:bj,part:complex} before meeting the filtration, MP--DXZ, and
degeneration interfaces in \Cref{part:comparisons}.

The exposition is addressed to graduate students.  We assume familiarity
with the language of schemes or compact Kähler manifolds, but no prior
knowledge of Berkovich geometry.  Definitions are followed quickly by
examples, and long proofs are divided into steps.  We cite the Stacks Project
for general scheme-theoretic facts whose proof would interrupt the main
argument.  We likewise cite \emph{Singularities in Global Pluripotential
Theory} for some deep complex-analytic foundations, while restating every
definition and interface needed here.

Several conventions deserve emphasis.  We use genuine nets whenever the
topology is not known to be sequential.  We do not subtract two infinite
values.  A zero Monge--Ampère mass is never treated as a harmless limiting
case without checking the definition.  Original theorems are first stated
under their original hypotheses; later improvements are stated separately.
This is especially important for residue characteristic, resolution of
singularities, continuity of envelopes, and regularity of Monge--Ampère
solutions.

The revised 2025 work of Chambert--Loir and Ducros is a principal mathematical
source for the local theory, but not the organizational model of
\cref{part:cld}.
There we proceed from explicit calculations on polyhedra, curves, and tori to
the invariant language.  Calibrations appear when integration requires them,
and abstract tropical spaces appear only after their role is visible.

Research questions are handled under a separate verification protocol.  An
old question already settled in the literature is reported as history, not
recycled as an open problem.  A false extension is replaced by a
counterexample.  A new argument is included only after a separate adversarial
audit.  No mathematical statement becomes true by being desirable.

This second draft retains repetition that helps readers cross between
formalisms; repeated proofs are replaced by precise references.
